\subsection{12 августа. Перевал Кашкасу.}
\textit{Метеоусловия: Солнечно, прохладно. }

\begin{figure}[h!]
	\centering
	\includegraphics[angle=0, width=0.7\linewidth]{../pics/maps/12_map.jpg}
	\label{fig:12_map}
	
\end{figure}


В 8:30 завтрак и выход. Группа решила перевыполнить план и закончить поход в назначенный день даже с учётом днёвки.  Идём активно. В 9:25	прошли грязевые топи.

\begin{figure}[h!]
	\centering
	\includegraphics[angle=0, width=0.7\linewidth]{../pics/days/12_1.jpg}
	\label{fig:12_1}
	\caption{ Грязевые топи рядом с МН8. }
\end{figure}

\begin{figure}[h!]
	\centering
	\includegraphics[angle=0, width=0.7\linewidth]{../pics/days/12_2.jpg}
	\label{fig:12_2}
	\caption{Вид в сторону МН, на в. Марс.}
\end{figure}

Далее следовали вереницей озёра, точнее, каждые полчаса. Высота последнего озера – 3888 м. Рельеф не представлял из себя проблем: слежавшиеся камни без курумника и сыпухи.

\begin{figure}[h!]
	\centering
	\includegraphics[angle=0, width=0.7\linewidth]{../pics/days/12_3.jpg}
	\label{fig:12_3}
	\caption{Примерный рельеф между озёрами.}
\end{figure}

Поднимаемся к перевалу, кругом существенно больше костей и черепов, животных, чем было на пути к предыдущим перевалам, в том числе нашли старый человеческий череп, о нем упоминалось в отчетах наших предшественников, собственно подъем был зловещий. Шли по турикам, которые ближе к перевалу пропали.
11:18, высота 3877 м.,	забрались выше, искали записку в неправильном месте. 
В 11:57 3890 поднялись на перевал Кашка-Суу. Перевальная записка в банке из-под протеина брутально соответствует духу перевала.

\begin{figure}[h!]
	\centering
	\includegraphics[angle=0, width=0.7\linewidth]{../pics/days/12_4.jpg}
	\label{fig:12_4}
	\caption{Решили соответствовать крутизне последнего перевала за этот поход.}
\end{figure}

\begin{figure}[h!]
	\centering
	\includegraphics[angle=0, width=0.7\linewidth]{../pics/days/12_5.jpg}
	\label{fig:12_5}
	\caption{Фотография группы без лиц, но с руководом и Кардиналом.}
\end{figure}

\begin{figure}[h!]
	\centering
	\includegraphics[angle=0, width=0.7\linewidth]{../pics/days/12_6.jpg}
	\label{fig:12_6}
	\caption{Фотография группы с лицами.}
\end{figure}

Фотографии, перевальные записки и начинаем спуск.  В 13:35 высота 3721, идём по курумнику, который постепенно перерастает в сыпуху. 

\begin{figure}[h!]
	\centering
	\includegraphics[angle=0, width=0.7\linewidth]{../pics/days/12_7.jpg}
	\label{fig:12_7}
	\caption{С двух сторон катятся камни и эффектно тают снежники.}
\end{figure}

\begin{figure}[h!]
	\centering
	\includegraphics[angle=0, width=0.7\linewidth]{../pics/days/12_8.jpg}
	\label{fig:12_8}
	\caption{Начало спуска. }
\end{figure}

Турики были хорошо видны на всём протяжении спуска, аж до самого места обеда. В 14:31 (высота 3400) курумник заканчивается и начинается движение по сыпухе 

\begin{figure}[h!]
	\centering
	\includegraphics[angle=0, width=0.7\linewidth]{../pics/days/12_9.jpg}
	\label{fig:12_9}
	\caption{Группа идёт по сыпухе. }
\end{figure}

\begin{figure}[h!]
	\centering
	\includegraphics[angle=0, width=0.7\linewidth]{../pics/days/12_10.jpg}
	\label{fig:12_10}
	\caption{После сыпухи начинается зелёный участок и идёт он вплоть до долины. }
\end{figure}

В 15:20	( 3325 м.) группа встала на обед. 

\begin{figure}[h!]
	\centering
	\includegraphics[angle=0, width=0.7\linewidth]{../pics/days/12_11.jpg}
	\label{fig:12_11}
	\caption{Место обеда с видом на спуск с перевала Кашка-Суу. }
\end{figure}

Во время обеда до нас доходит группа Алексея Остапива и встают тут на место ночёвки. Наша группа, дабы силы ещё есть и идти в последний день хочется как можно меньше, принимает решение после обеда продолжить движение вниз к долине к МН10. В 17:38 продолжаем движение. Спуск не представляет сложности. Он приятный и простой, одна сплошная трава с протоптанными тропинками.

\begin{figure}[h!]
	\centering
	\includegraphics[angle=0, width=0.7\linewidth]{../pics/days/12_12.jpg}
	\label{fig:12_12}
	\caption{Вид на спуск с перевала Кашка-Суу. }
\end{figure}

В 18:25 спустились в долину и разбили лагерь. Состояние у группы хорошее, самочувствие у всех в норме. Высота 3000.

\begin{figure}[h!]
	\centering
	\includegraphics[angle=0, width=0.7\linewidth]{../pics/days/12_13.jpg}
	\label{fig:12_13}
	\caption{Группа идёт в долину к МН10 (выполаживание на берегу). }
\end{figure}

ЧХВ:	5:43
ГХВ:	7:37
\clearpage